\begin{tikzpicture}
\begin{semilogxaxis}[
    title={Diagrama de Bode de Magnitud (Valores Reales)},
    xlabel={Frecuencia $\omega$ (rad/s)},
    ylabel={Magnitud (dB)},
    grid=both,
    grid style={line width=.1pt, draw=gray!30},
    major grid style={line width=.2pt, draw=gray!60},
    xmin=1e-1, xmax=1e5,
    ymin=-50, ymax=30,
    width=0.9\columnwidth,
    height=8cm,
    legend style={
    nodes={scale=0.7, transform shape}
}
]

\addplot[thick, blue] coordinates {
    (0.1, -26.38)
    (1, -6.38)
    (8.333, 12.04)
    (166.667, 12.04)
    (1e5, -43.52)
};
\addlegendentry{Teórico (Asintótico)}

\addplot[mark=*, mark size=1.5pt, red, forget plot] coordinates {(8.333, 12.04)};
\node[anchor=south east, font=\small, text=black] at (axis cs:8.333, 12.04) {$\omega_{p1} = 8.3$};

\addplot[mark=*, mark size=1.5pt, red, forget plot] coordinates {(166.667, 12.04)};
\node[anchor=south west, font=\small, text=black] at (axis cs:166.667, 12.04) {$\omega_{p2} = 166.7$};

% Datos de LTSpice
\addplot[thick, dashed, orange] table [
    skip first n=1,
    x expr=\thisrowno{0} * 2 * pi,
    y index=3
] {sources/circuito3_cleaned.txt};
\addlegendentry{Simulación LTSpice}

\end{semilogxaxis}
\end{tikzpicture}
